В работе была исследована РСМО с потерями вида $GI/G/N/N/K$, в которой каждая поступающая заявки занимает 1 прибор и случайный объем ограниченного ресурса. Основное внимание было уделено анализу чувствительности стационарных характеристик системы к форме распределения входящего потока и распределения объема работы заявки при фиксированных первых моментах.

Начиная с классической модели Эрланга $M/M/N/N$ рассматривались все более общие случаи. Сначала при генерализации распределения времени работы в $M/G/N/N$ системе проявляется свойство инвариантности к режиму работы заявки. Затем рассмотрелись ресурсные системы $M/G/N/N/R$, которые добавили дополнительное условие к причинам отказа заявки. И наконец с учетом современных исследований трафика было генерализировано распределение входящего потока заявок. Наличие случайного ресурсного требования и отказ от пуассоновского входящего потока приводят к потере марковского свойства и резкому росту пространства состояний.

Во второй главе была формализована математическая модель РСМО с потерями, учитывающая ограничения по прибору ($N$) и ресурсу ($R$). Состояние системы в произвольный момент времени описано через текущее число заявок $\xi(t)$ и суммарный объем занятого ими ресурса $\delta(t)$. Поступающая заявка с требованием $r_i$ теряется при выполнении хотя бы одного из условий: $\xi(t-) = N$ или $\delta(t-) + r_i > R$. В качестве целевых метрик эффективности исследованы вероятность отказа ($P_{\mathrm{loss}}$), абсолютная пропускная способность ($\theta$), математическое ожидание занятого ресурса $\mathbb{E}[\delta]$ и оценка стационарного распределения вероятностей состояний $\pi(k)$. 

Для оценки стационарных характеристик системы был разработан алгоритм и дискретно-событийного имитационного моделирования методом Монте-Карло. Чтобы отделить влияние средней нагрузки от влияния формы распределений, в вычислительном эксперименте задавались семейства распределений входящего потока и объема работы заявки с идентичными математическими ожиданиями, но различными коэффициентами вариации $c_v$. 

Для распределения интервалов между заявками были выбраны базовое пуассоновское распределение ($c_v = 1$), распределения Эрланга 4-го и 2-го порядков ($c_v = 0.5$ и $c_v \approx 0.71$) соответственно, моделирующие регулярные потоки и гиперэкспоненциальное распределение ($c_v > 1$, моделирующее пачечный характер трафика и тяжелые хвосты).

Для распределения времени работы заявки было взять базовое детерменированное распределение ($c_v = 0$), экспоненциальное ($c_v = 1$), распределение Эрланга 4го порядка ($c_v = 0.5$) и гиперэкспоненциальное ($c_v > 1$). 

Особое внимание в работе уделено архитектуре программного комплекса. Учитывая отсутствие буфера ожидания и, как следствие, отсутствие необходимости динамически выделять память под отслеживание объектов в очереди, каждая репликация Монте-Карло представляет собой независимый процесс. Все $960$ потоков, полученные благодаря трем значениям $\lambda$, четырем распределениям интенсивности входящего потока, четырем распределениям времени работы заявки при 40 репликацияз  запускаются параллельно на ядрах CUDA. Это обеспечивает прирост производительности по сравнению с CPU. 

Подтвердилось свойство инвариантности ресурсной системы с потерями к распределению времени обслуживания при пуассоновском входящем потоке. При сохранении одинакового среднего времени обслуживания переход от детерминированного распределения к экспоненциальному, эрланговскому или гиперэкспоненциальному распределению практически не изменял вероятность отказа, пропускную способность и среднюю занятость системы. Наблюдаемые расхождения находились на уровне статистического шума метода Монте-Карло.

Тип входящего потока оказался более существенным фактором чувствительности РСМО. При одинаковой средней интенсивности $\lambda$ регулярные потоки Эрланга приводили к снижению вероятности отказа и более устойчивому заполнению системы, тогда как гиперэкспоненциальный входящий поток, моделирующий всплески поступлений, существенно увеличивал вероятность отказа. 

Различная интенсивность входящего потока $\lambda$, моделирующая низкую, среднюю и высокую нагрузку, оказалась самым существенным фактором влияния на вероятность отказа, пропускную способность и среднюю занятость системы. 

Совместный анализ распределений входящего потока и времени обслуживания показал, что при непуассоновском входе свойство инвариантности нарушается. Для регулярных входящих потоков наилучшие результаты дает более регулярное обслуживание, тогда как при всплесковом гиперэкспоненциальном входе повышение вариативности времени обслуживания может частично снижать вероятность отказа за счет более неравномерного освобождения ресурсов. Таким образом, распределение времени обслуживания само по себе не имеет универсального положительного или отрицательного эффекта: его влияние проявляется через взаимодействие с характером входящего потока.

Таким образом в ходе работы была построена и исследована имитационная модель ресурсной системы с потерями $GI/G/N/N/R$, разработан программно-вычислительный комплекс для проведения серии экспериментов, а также выявлены основные факторы чувствительности стационарных характеристик. Полученные результаты показывают, что для ресурсных СМО с потерями в первую очередь необходимо учитывать среднюю нагрузку и тип входящего потока, затем --- взаимодействие входящего потока с распределением времени обслуживания. 



